\documentclass[14pt,a4paper]{extarticle}

\usepackage[T2A]{fontenc}
\usepackage[utf8]{inputenc}
\usepackage[russian]{babel}
\usepackage{geometry}
\usepackage{xcolor}
\usepackage{graphicx}
\usepackage{enumitem}
\usepackage{titlesec}
\usepackage{parskip}

\geometry{left=18mm,right=18mm,top=16mm,bottom=18mm}
\setlength{\parindent}{0pt}
\setlength{\parskip}{0.45em}
\linespread{1.12}
\pagestyle{plain}
\emergencystretch=2em
\sloppy

\definecolor{accent}{RGB}{208,229,205}
\definecolor{accentline}{RGB}{148,180,144}
\definecolor{textgray}{RGB}{70,70,70}

\titleformat{\section}
  {\Large\bfseries}
  {}
  {0pt}
  {}

\newcommand{\Slide}[2]{%
  \clearpage
  \section*{Слайд #1. #2}
  \noindent\colorbox{accent}{\parbox{\dimexpr\linewidth-2\fboxsep\relax}{\textbf{Текст для выступления}}}
  \vspace{0.8em}
}

\newcommand{\SlidePreview}[1]{%
  \begin{center}
  \fbox{\includegraphics[page=#1,width=0.88\linewidth]{output/presentation.pdf}}
  \end{center}
  \vspace{0.6em}
}

\begin{document}

\Slide{1}{Титульный слайд}
\SlidePreview{1}
Добрый день. Тема моей работы --- <<Разработка программной системы фиксации городских событий>>.

Данная работа выполнена в рамках преддипломной практики и объединяет результаты двух семестров. В первом семестре были реализованы серверная часть и диспетчерская веб-панель, во втором --- мобильное приложение для жителей. В настоящей работе представлена итоговая программная система, включающая все три компонента.

\Slide{2}{Анализ текущей проблемы}
\SlidePreview{2}
{\small
\setlength{\parskip}{0.2em}
В настоящее время в городе Кирове отсутствует единый цифровой канал для сообщения о проблемах городской инфраструктуры. Информация о неисправностях поступает в муниципальные службы по нескольким разрозненным каналам:

а) по телефону;

б) по электронной почте;

в) через личные сообщения.

Такой подход создаёт ряд организационных затруднений: обращения могут теряться, их сложно распределять между исполнителями, а контроль обработки в реальном времени невозможен.

Не менее существенные проблемы испытывают жители: отсутствует возможность оперативно зафиксировать проблему с мобильного устройства, житель не видит уже созданные обращения поблизости и не может отследить статус обработки.

Без интерактивной карты и централизованного учёта диспетчер не видит общую картину проблем города и вынужден тратить дополнительное время на ручную маршрутизацию обращений.
}

\Slide{3}{Цель и задачи работы}
\SlidePreview{3}
Цель работы --- разработать программную систему для фиксации и обработки обращений о проблемах городской инфраструктуры города Кирова, обеспечивающую полный цикл обработки обращений --- от подачи жителем через мобильное приложение до контроля исполнения диспетчером через веб-панель.

Система включает три компонента: серверную часть с базой данных PostgreSQL и расширением PostGIS, диспетчерскую веб-панель с интерактивной картой на базе Яндекс.Карт и кроссплатформенное мобильное приложение на React Native для жителей.

Задачи работы: провести анализ предметной области, спроектировать структуру системы и базу данных, реализовать серверную часть с программным интерфейсом, разработать веб-панель диспетчера, создать мобильное приложение для жителей и провести тестирование всех компонентов.

\Slide{4}{Анализ существующих решений}
\SlidePreview{4}
Были рассмотрены существующие аналоги: платформа обратной связи <<Госуслуги. Решаем вместе>>, <<Наш город>> (Москва), <<Добродел>> (Московская область) и FixMyStreet (Великобритания).

Проведённый анализ показал, что среди рассмотренных аналогов отсутствует решение, применимое к городу Кирову и обеспечивающее полный набор функций: автоматическое определение геолокации, отслеживание статуса, просмотр обращений на карте и механизм дедупликации по радиусу.

Поэтому был сделан вывод о целесообразности разработки собственной системы.

\Slide{5}{Технологический стек}
\SlidePreview{5}
Для реализации системы был выбран следующий технологический стек.

Серверная часть построена на ASP.NET Core с использованием языка C\#. PostgreSQL выбрана в первую очередь потому, что используется вместе с PostGIS для работы с координатами и геопространственными запросами. Функция ST\_DWithin позволяет эффективно искать обращения в заданном радиусе.

Диспетчерская веб-панель реализована на React с TypeScript и Яндекс.Картами для интерактивной визуализации обращений.

Мобильное приложение --- React Native с инструментарием Expo. Главный довод --- кроссплатформенность и преемственность с веб-панелью. Expo предоставляет готовые модули для камеры и геолокации.

\Slide{6}{Структура системы}
\SlidePreview{6}
Здесь представлена общая структура программной системы.

Мобильное приложение и веб-панель диспетчера --- это два клиента, которые обращаются к одному серверу на C\#. Сервер работает с PostgreSQL и PostGIS.

Таким образом, оба клиента используют единый программный интерфейс, что обеспечивает согласованность данных.

\Slide{7}{Особенности проектирования базы данных}
\SlidePreview{7}
Основной сущностью базы данных является обращение.

В обращении хранятся основные данные о проблеме: описание, координаты, статус, дата создания и дата обновления.

Также обращение связано с категорией проблемы, пользователем, который его создал, медиафайлами, историей статусов и ответственными службами.

Именно такая структура позволяет не просто хранить заявку, а полноценно сопровождать её на всех этапах обработки.

\Slide{8}{Реализация: дашборд диспетчера}
\SlidePreview{8}
На данном слайде представлен дашборд веб-панели диспетчера.

Этот экран содержит сводные показатели по обращениям, графики распределения по категориям и статусам, а также список последних поступивших записей. Такой интерфейс позволяет быстро оценить текущую ситуацию и перейти к обработке конкретного обращения.

Таким образом, дашборд выполняет функцию стартовой аналитической панели для диспетчера.

\Slide{9}{Реализация: интерактивная карта диспетчера}
\SlidePreview{9}
Центральным рабочим экраном диспетчера является интерактивная карта обращений на базе Яндекс.Карт.

На карте отображаются маркеры обращений, причём цвет маркера зависит от текущего статуса. Диспетчер может использовать фильтрацию и быстро находить проблемные точки на городской территории.

Такой подход повышает наглядность обработки обращений и позволяет оценивать ситуацию не только списком, но и пространственно.

\Slide{10}{Реализация: список обращений}
\SlidePreview{10}
На этом слайде показана страница со списком обращений.

Она содержит табличный перечень записей с основными полями, вкладки по статусам и средства быстрого поиска. Такой формат удобен для последовательной обработки заявок и позволяет быстро переходить к детальной карточке нужного обращения.

По сути, это основной режим работы со структурированным перечнем поступивших сообщений.

\Slide{11}{Реализация: карточка обращения}
\SlidePreview{11}
Карточка обращения содержит полную информацию по конкретной заявке.

Здесь отображаются описание проблемы, категория, текущий статус, история обработки, мини-карта с местоположением и данные ответственной службы. Также диспетчеру доступна форма смены статуса с добавлением комментария.

Именно этот экран обеспечивает непосредственную диспетчерскую обработку обращения и фиксацию результата.

\Slide{12}{Мобильное приложение: создание обращения}
\SlidePreview{12}
Это ключевой экран мобильного приложения. Житель выбирает категорию проблемы, вводит описание, фотографирует проблему и указывает местоположение на карте --- по GPS или вручную.

При обнаружении аналогичного обращения в радиусе 50~метров приложение предлагает подтвердить существующее обращение вместо создания нового. Этот механизм дедупликации реализован на стороне сервера.

\Slide{13}{Мобильное приложение: карта обращений}
\SlidePreview{13}
На карте мобильного приложения отображаются маркеры обращений, цвет которых зависит от статуса: новое, принято, в работе, решено.

Реализованы фильтры по категориям, кнопка определения текущего местоположения и плавающая кнопка для быстрого перехода к созданию обращения.

\Slide{14}{Мобильное приложение: список и карточка обращений}
\SlidePreview{14}
Пользователь видит свои обращения в виде списка. Есть вкладки --- все, активные и завершённые.

При нажатии на обращение открывается полная карточка --- фото, адрес, описание и вся история изменений статуса. Другие жители могут нажать <<Тоже вижу проблему>>, и это увеличит приоритет обращения.

\Slide{15}{Тестирование}
\SlidePreview{15}
Проведено тестирование всех компонентов системы.

Автоматизированные тесты серверной части выполнены на платформе xUnit --- модульные и интеграционные. Модульные тесты мобильного приложения реализованы на Vitest --- 20 тестовых файлов. Также выполнено тестирование REST API через Swagger~UI и функциональное тестирование веб-панели диспетчера.

Все тесты завершились успешно.

\Slide{16}{Перспективы развития}
\SlidePreview{16}
Дальнейшее развитие системы может включать внедрение push-уведомлений на уровне операционной системы, расширение возможностей работы при нестабильном сетевом соединении, развитие модуля городских событий.

Кроме того, перспективным направлением является интеграция с внешними картографическими сервисами для обратного геокодирования и реализация ролевой модели с несколькими уровнями доступа для диспетчеров различных служб.

\Slide{17}{Заключение}
\SlidePreview{17}
В результате выполнения работы спроектирована и реализована программная система фиксации городских событий, предназначенная для жителей города Кирова и диспетчеров муниципальных служб.

Система включает серверную часть на ASP.NET Core с PostgreSQL и PostGIS, диспетчерскую веб-панель на React с Яндекс.Картами и кроссплатформенное мобильное приложение на React Native с Expo.

Все задачи работы выполнены в полном объёме. Разработанная программная система обеспечивает полный цикл обработки обращений --- от подачи жителем до контроля диспетчером.

Спасибо за внимание. Готова ответить на вопросы.

\end{document}
